% =====================================================================
% CHƯƠNG 5: KẾT QUẢ VÀ ĐÁNH GIÁ
% ĐỀ TÀI: XÂY DỰNG HỆ THỐNG ĐỌC NHÃN MÁC VÀ CẢNH BÁO DỊ ỨNG THỰC PHẨM
% =====================================================================

\chapter{Kết quả và đánh giá}
\label{ch:ketqua}

\section{Kết quả đạt được}

\subsection{Kết quả xây dựng hệ thống}

Hệ thống Allergy Warning được xây dựng nhằm hỗ trợ người dùng đọc
thông tin trên nhãn thực phẩm và nhận biết các thành phần liên quan
đến dị ứng. Các chức năng chính gồm:

\begin{itemize}
	\item Đăng ký, đăng nhập và quản lý hồ sơ cá nhân.
	\item Thiết lập hồ sơ dị ứng và lựa chọn các dị nguyên quan tâm.
	\item Tải ảnh nhãn thực phẩm để thực hiện phân tích.
	\item Sử dụng Gemini API để nhận diện và trích xuất thông tin trên nhãn.
	\item Chuẩn hóa thành phần, đối chiếu với dữ liệu dị nguyên và hồ sơ người dùng.
	\item Hiển thị cảnh báo, lưu và tra cứu lịch sử quét.
\end{itemize}

Hệ thống sử dụng Flask làm máy chủ, TiDB Cloud lưu trữ dữ liệu,
Cloudinary lưu trữ hình ảnh và Gemini API hỗ trợ phân tích nhãn.
Open Food Facts được sử dụng làm nguồn dữ liệu tham khảo; ứng dụng
được triển khai trên Render.

\subsection{Kết quả phân tích và cảnh báo}

Quy trình phân tích gồm tiếp nhận ảnh, lưu ảnh trên Cloudinary,
gửi ảnh đến Gemini API, chuẩn hóa thông tin và đối chiếu thành phần
với dữ liệu dị nguyên. Kết quả được tổng hợp để hiển thị cho người dùng.

Hệ thống phân biệt thành phần được liệt kê trên nhãn với cảnh báo
về khả năng chứa vết dị nguyên. Đồng thời, hệ thống thể hiện riêng
dị nguyên được nhận diện và dị nguyên trùng với hồ sơ cá nhân.
Kết quả chỉ có tính chất tham khảo, không khẳng định sản phẩm an toàn
tuyệt đối.

\section{Đánh giá các chức năng}

Việc đánh giá tập trung vào khả năng hoạt động và tính phù hợp của
kết quả ở các chức năng chính.

\subsection{Quản lý tài khoản và hồ sơ dị ứng}

Kiểm tra khả năng đăng ký, đăng nhập, cập nhật hồ sơ và lưu lựa chọn
dị nguyên. Dữ liệu hồ sơ và lịch sử quét cần được truy xuất đúng theo
tài khoản đang đăng nhập.

\subsection{Phân tích hình ảnh và nhận diện nội dung}

Đối chiếu thông tin Gemini API trả về với nhãn gốc. Ảnh rõ, đủ sáng
giúp quá trình nhận diện thuận lợi hơn; ảnh mờ, nghiêng hoặc có chữ
nhỏ có thể làm thiếu hoặc sai thông tin.

\subsection{Đối chiếu dị nguyên}

Kiểm tra khả năng chuẩn hóa tên thành phần và phát hiện dị nguyên.
Cần chú ý các tên gọi tương đương, những thành phần gần giống nhau
và trường hợp dị nguyên được nhận diện nhưng không nằm trong hồ sơ
dị ứng của người dùng.

\subsection{Hiển thị kết quả và lịch sử quét}

Kiểm tra việc hiển thị thông tin sản phẩm, thành phần và cảnh báo.
Khi phân tích thất bại hoặc thiếu dữ liệu, hệ thống cần thông báo rõ.
Lịch sử phải được lưu và truy xuất đúng tài khoản.

\section{Phương pháp và tiêu chí đánh giá}

\subsection{Phương pháp đánh giá}

Sử dụng ảnh nhãn thực phẩm làm dữ liệu kiểm thử. Thông tin thành phần
và cảnh báo trên nhãn gốc được ghi nhận để so sánh với kết quả hệ thống.
Tập ảnh nên có nhiều tình huống như ảnh rõ, ảnh chữ nhỏ, ảnh nhiều
thành phần và ảnh có cảnh báo về khả năng chứa vết dị nguyên.

\subsection{Các tiêu chí đánh giá}

Hệ thống được đánh giá dựa trên các tiêu chí liên quan đến khả năng
phân tích nhãn thực phẩm, nhận diện thành phần và cảnh báo dị ứng.
Cụ thể gồm:

\begin{itemize}
	\item \textbf{Khả năng trích xuất thông tin:}
	Đánh giá mức độ đầy đủ và chính xác của thông tin được trích xuất
	từ hình ảnh nhãn thực phẩm, bao gồm tên sản phẩm, thương hiệu,
	thành phần và thông tin cảnh báo.
	
	\item \textbf{Khả năng nhận diện thành phần:}
	Kiểm tra khả năng đọc và chuẩn hóa danh sách thành phần thực phẩm,
	kể cả khi thông tin trên nhãn có cách viết hoặc cách diễn đạt khác nhau.
	
	\item \textbf{Khả năng phát hiện chất gây dị ứng:}
	Đánh giá khả năng đối chiếu thành phần sản phẩm với danh mục
	chất gây dị ứng và hồ sơ dị ứng của người dùng để xác định
	các thành phần có nguy cơ gây dị ứng.
	
	\item \textbf{Khả năng đưa ra cảnh báo:}
	Kiểm tra việc hiển thị đầy đủ các chất gây dị ứng được phát hiện,
	đồng thời phân biệt thành phần có trong sản phẩm với cảnh báo
	về khả năng chứa dấu vết chất gây dị ứng.
	
	\item \textbf{Khả năng xử lý các trường hợp đầu vào:}
	Đánh giá khả năng xử lý hình ảnh rõ nét, hình ảnh mờ,
	thiếu sáng, chứa nhiều thành phần hoặc không đủ thông tin
	để phân tích.
	
	\item \textbf{Tính ổn định của hệ thống:}
	Kiểm tra khả năng xử lý khi dịch vụ Gemini API gặp lỗi,
	hình ảnh không tải lên được hoặc không thể truy xuất thông tin
	sản phẩm, đồng thời đảm bảo hệ thống thông báo phù hợp
	cho người dùng.
\end{itemize}

Các tiêu chí trên được sử dụng để xem xét mức độ đáp ứng yêu cầu
của hệ thống. Kết quả đánh giá là cơ sở để xác định những ưu điểm,
hạn chế và khả năng cải thiện trong quá trình phát triển ứng dụng.
\subsection{Đánh giá định lượng}

Khi có dữ liệu kiểm thử được đối chiếu, có thể sử dụng Precision,
Recall và F1-score để đánh giá khả năng phát hiện dị nguyên:

\[
Precision = \frac{TP}{TP + FP}
\]

\[
Recall = \frac{TP}{TP + FN}
\]

\[
F1\text{-}score =
2 \times \frac{Precision \times Recall}{Precision + Recall}
\]

Trong đó, TP là số trường hợp phát hiện đúng, FP là số trường hợp
cảnh báo sai và FN là số trường hợp bỏ sót. Các chỉ số cần được tính
từ kết quả kiểm thử thực tế.

\begin{table}[htbp]
	\centering
	\caption{Tổng hợp kết quả đánh giá định lượng}
	\label{tab:ch5-ketqua-dinhluong}
	\small
	\renewcommand{\arraystretch}{1.15}
	\setlength{\tabcolsep}{4pt}
	\begin{tabularx}{\textwidth}{|>{\raggedright\arraybackslash}p{4.2cm}|>{\centering\arraybackslash}X|}
		\hline
		\textbf{Chỉ số} & \textbf{Kết quả} \\ \hline
		Tổng số ảnh kiểm thử & \textit{Điền số liệu thực tế} \\ \hline
		TP, FP, FN & \textit{Điền số liệu thực tế} \\ \hline
		Precision, Recall, F1-score & \textit{Tính từ dữ liệu kiểm thử} \\ \hline
	\end{tabularx}
\end{table}

Nếu chưa có tập dữ liệu đủ tin cậy, cần trình bày đánh giá định tính
và không tự ước lượng các chỉ số.

\section{Bàn luận}

Kết quả hệ thống chịu ảnh hưởng bởi ba yếu tố chính. Thứ nhất,
chất lượng ảnh không tốt có thể khiến OCR và Gemini API nhận diện
thiếu hoặc sai nội dung. Thứ hai, dữ liệu dị nguyên và quy tắc chuẩn
hóa chưa đầy đủ có thể dẫn đến bỏ sót hoặc đối chiếu nhầm thành phần.
Thứ ba, Gemini API và các dịch vụ bên ngoài có thể gặp lỗi kết nối
hoặc trả về dữ liệu không phù hợp.

Nhìn chung, Allergy Warning giúp người dùng tiếp cận thông tin thành
phần và cảnh báo thuận tiện hơn thông qua việc phân tích ảnh và đối
chiếu với hồ sơ dị ứng. Tuy nhiên, hệ thống chỉ hỗ trợ tham khảo,
không thay thế việc kiểm tra nhãn gốc hoặc tư vấn chuyên môn.

\section{Hạn chế}

Hệ thống còn một số hạn chế:

\begin{itemize}
	\item Kết quả nhận diện phụ thuộc vào chất lượng hình ảnh và phản hồi của Gemini API.
	\item Dữ liệu dị nguyên và các biến thể tên thành phần chưa bao quát mọi trường hợp.
	\item Quy tắc chuẩn hóa có thể bỏ sót hoặc đối chiếu nhầm các tên gọi gần giống.
	\item Việc đánh giá còn phụ thuộc vào số lượng và độ đa dạng của ảnh kiểm thử.
\end{itemize}

\section{Hướng cải thiện}

Trong thời gian tới, hệ thống có thể được cải thiện bằng cách:

\begin{itemize}
	\item Bổ sung kiểm tra chất lượng ảnh và hướng dẫn người dùng chụp ảnh phù hợp.
	\item Mở rộng, kiểm chứng dữ liệu dị nguyên và tên gọi tương đương.
	\item Hoàn thiện quy tắc đối chiếu để hạn chế cảnh báo sai.
	\item Tăng cường xử lý lỗi khi Cloudinary, Gemini API hoặc cơ sở dữ liệu gặp sự cố.
	\item Xây dựng tập ảnh kiểm thử đa dạng để đánh giá chính xác hơn.
\end{itemize}

\section{Kết luận chương}

Chương 5 đã trình bày kết quả xây dựng hệ thống Allergy Warning,
các chức năng chính và phương pháp đánh giá. Qua đó, chương cũng
chỉ ra những yếu tố ảnh hưởng đến kết quả phân tích, các hạn chế
hiện tại và hướng cải thiện trong tương lai.
